We introduce constituent-representative geographic distance as an outcome measure for
evaluating congressional redistricting plans. For each of the 435 congressional districts
in the 2022 enacted maps, we compute the mean driving time from each census tract's
population-weighted centroid to the representative's primary district office address,
as listed on \texttt{house.gov}. Routing is performed via the OSRM (Open Source Routing
Machine) API over the full OpenStreetMap road network. We document a correlation of
$r = -0.67$ (95\% CI: $[-0.72, -0.61]$) between district Polsby-Popper compactness and
mean constituent driving time across all 435 districts: less compact districts impose
longer travel times on their constituents.

Comparing enacted maps to \textsc{bisect} algorithmic plans computed from the same
census data, we find that \textsc{bisect} districts achieve a mean constituent driving
time of 34 minutes vs.\ 48 minutes for enacted gerrymanders in the 12 states with
documented partisan gerrymandering (Efficiency Gap $|EG| > 0.08$). The gap is driven
by the ``neck'' mechanism: gerrymandered districts connect geographically distant
communities through narrow corridors, inflating the perimeter-to-area ratio (reducing
Polsby-Popper) while simultaneously stretching the maximum distance between district
constituents. We characterize this mechanism geometrically and show that the correlation
between district ``neck width'' (minimum corridor width as a fraction of district diameter)
and mean constituent driving time is $r = 0.71$, stronger than the Polsby-Popper
correlation.

The result has practical implications for redistricting litigation: constituent travel
distance is a direct, welfare-relevant harm that is predictable from map geometry without
requiring election data. Courts evaluating compact redistricting criteria can cite
constituent access as a concrete, non-partisan outcome that compactness promotes.
